\section{Conclusion}
\label{sec:conclusion}

This paper addressed a fundamental architectural gap in agricultural artificial intelligence: the absence of a horizontal, local-first agent platform that can adapt to diverse crop-protection domains through declarative configuration rather than vertical re-engineering. Against this backdrop, we established three research objectives: (1) verifying MCP-based decoupling between a Rust agent kernel and a local agricultural knowledge base; (2) proposing and validating a Domain-as-Config methodology for zero-code adaptation; and (3) evaluating the resulting system on offline knowledge-retrieval and diagnostic tasks.

All three objectives were met. The integration tests demonstrate that Clarity and devbase can communicate reliably over MCP stdio transport, including robust handling of Content-Length-framed JSON-RPC messages. The benchmark experiments confirm that loading an agricultural `DomainConfig` at runtime elevates diagnostic retrieval accuracy from 20\% (uninformed random sampling) to 100\% (perfect Hit@3), with a KeywordRecall@1 of 99.05\%, while maintaining an end-to-end latency of approximately 16--19~ms. Configuration migration---switching from a generic software-engineering profile to the agricultural profile---requires less than a millisecond and involves zero changes to the Rust source code.

The principal contributions of the work are therefore threefold. First, we presented the first MCP-based, configuration-driven agent architecture specifically designed for agricultural knowledge management. Second, we implemented and open-sourced a standards-compliant Rust MCP client/server stack with dynamic tool discovery and resilient stdio transport. Third, we provided empirical evidence that cross-crop agricultural adaptation can be reduced from hours of coding to minutes of configuration editing, without sacrificing retrieval accuracy or interactive latency.

Looking ahead, we believe that the combination of local-first runtime architectures, standardized tool protocols such as MCP, and declarative domain configuration represents a promising path toward equitable, privacy-preserving, and cost-effective agricultural AI. Future efforts will focus on scaling the knowledge base, integrating on-device vision models as additional MCP tools, and deploying the full conversational pipeline with local LLM inference in real rural field settings.
